\documentclass[11pt]{article}
\usepackage{amsmath,amssymb,amsthm}
\usepackage{algorithm}
\usepackage{algorithmic}
\usepackage{fullpage}
\usepackage{hyperref}
\newtheorem{theorem}{Theorem}
\newtheorem{lemma}{Lemma}
\newtheorem{assumption}{Assumption}
\newcommand{\E}{\mathbb{E}}
\newcommand{\R}{\mathbb{R}}
\newcommand{\norm}[1]{\left\|#1\right\|}
\newcommand{\block}[1]{\mathcal{B}_{#1}}
\newcommand{\grad}[1]{\nabla_{#1} #1}
\begin{document}

\section*{Algorithm Name}
\textbf{Randomized Block Coordinate Descent (R‑BCD) for Non‑Convex Smooth Objectives}

\section*{Mathematical Formulation}
Let $f:\R^{d}\rightarrow\R$ be a differentiable (possibly non‑convex) objective.  
Partition the coordinate set $\{1,\dots ,d\}$ into $B$ disjoint blocks 
\[
\{1,\dots ,d\}= \bigcup_{b=1}^{B}\block{b},\qquad 
\block{b}\cap\block{b'}=\emptyset\;(b\neq b').
\]
For a vector $w\in\R^{d}$ we denote by $w_{\block{b}}$ the sub‑vector consisting of the coordinates in block $\block{b}$ and by $\nabla_{\block{b}} f(w)$ the corresponding partial gradient.

At iteration $t=0,1,2,\dots$:

\begin{enumerate}
\item Sample a block $i_{t}\in\{1,\dots ,B\}$ uniformly at random:
\[
\Pr(i_{t}=b)=\frac{1}{B},\qquad \forall b.
\]
\item Compute the block gradient $\nabla_{\block{i_t}} f(w_t)$.
\item Update only the selected block with a (possibly time‑varying) stepsize $\eta_t>0$:
\[
w_{t+1}= w_t - \eta_t \, U_{i_t}\,\nabla_{\block{i_t}} f(w_t),
\]
where $U_{i_t}\in\R^{d\times |\block{i_t}|}$ is the zero‑padding matrix that maps a block vector to the full space (i.e. $(U_{i_t}v)_{j}=v_{j}$ if $j\in\block{i_t}$ and $0$ otherwise).
\end{enumerate}

\section*{Pseudocode}
\begin{algorithm}[H]
\caption{Randomized Block Coordinate Descent (R‑BCD)}
\begin{algorithmic}[1]
\STATE \textbf{Input:} initial point $w_{0}\in\R^{d}$, stepsize schedule $\{\eta_t\}_{t\ge 0}$, number of blocks $B$.
\FOR{$t=0,1,2,\dots,T-1$}
    \STATE Sample block $i_t\sim\text{Uniform}\{1,\dots ,B\}$ \hfill\COMMENT{Random block}
    \STATE Compute block gradient $g_t:=\nabla_{\block{i_t}} f(w_t)$
    \STATE $w_{t+1}= w_t - \eta_t\,U_{i_t} g_t$ \hfill\COMMENT{Update selected block}
\ENDFOR
\STATE \textbf{Return} $w_T$ (or the averaged iterate $\bar w_T=\frac{1}{T}\sum_{t=0}^{T-1}w_t$)
\end{algorithmic}
\end{algorithm}

\section*{Dry Run Example}
Consider the univariate function $f(w)=(w-3)^2$ (hence $d=1$, $B=1$).  
The full gradient equals $\nabla f(w)=2(w-3)$.  
Take $w_0=0$, constant stepsize $\eta_t=\eta=0.1$, and run three iterations.

\begin{center}
\begin{tabular}{c|c|c|c}
$t$ & $w_t$ & $\nabla f(w_t)$ & $f(w_t)$\\ \hline
0 & $0$ & $2(0-3)=-6$ & $(0-3)^2=9$\\
1 & $w_1=0-0.1(-6)=0.6$ & $2(0.6-3)=-4.8$ & $(0.6-3)^2=5.76$\\
2 & $w_2=0.6-0.1(-4.8)=1.08$ & $2(1.08-3)=-3.84$ & $(1.08-3)^2=3.6864$\\
3 & $w_3=1.08-0.1(-3.84)=1.464$ & $2(1.464-3)=-3.072$ & $(1.464-3)^2=2.3660$
\end{tabular}
\end{center}
The objective value decreases at each step, illustrating the behaviour of the algorithm in the simplest (single‑block) setting.

\newpage
\section*{Assumptions}
\begin{assumption}[Blockwise $L$‑smoothness]\label{ass:smooth}
For every block $b\in\{1,\dots ,B\}$ there exists a constant $L_b>0$ such that for all $w,\tilde w\in\R^{d}$,
\[
f\big(w+U_{b}(v-\nabla_{\block{b}} f(w))\big)
\le 
f(w)+\langle \nabla_{\block{b}} f(w),v\rangle+\frac{L_b}{2}\|v\|^{2},
\quad\forall v\in\R^{|\block{b}|}.
\]
Equivalently, the block gradient is $L_b$‑Lipschitz:
\[
\norm{\nabla_{\block{b}} f(w)-\nabla_{\block{b}} f(\tilde w)}\le L_b \norm{U_{b}^{\top}(w-\tilde w)}.
\]
\end{assumption}

\begin{assumption}[Bounded below]\label{ass:lower}
The objective is bounded from below:
\[
f^{\star}:=\inf_{w\in\R^{d}} f(w)>-\infty .
\]
\end{assumption}

\begin{assumption}[Stepsize]\label{ass:stepsize}
The stepsizes satisfy
\[
\eta_t = \frac{\eta}{\sqrt{t+1}},\qquad \eta>0,
\]
or more generally $\eta_t\le \frac{1}{L_{\max}}$ where $L_{\max}:=\max_{b}L_b$.
\end{assumption}

\section*{Main Convergence Theorem}
\begin{theorem}[Expected Gradient Norm Decay]\label{thm:main}
Under Assumptions \ref{ass:smooth}–\ref{ass:stepsize}, let $\{w_t\}_{t\ge0}$ be the iterates generated by R‑BCD. Define the average iterate
\[
\bar w_T:=\frac{1}{T}\sum_{t=0}^{T-1} w_t .
\]
Then for any $T\ge 1$,
\[
\frac{1}{T}\sum_{t=0}^{T-1}\E\!\big[\norm{\nabla f(w_t)}^{2}\big]
\le
\frac{2B\big(f(w_0)-f^{\star}\big)}{\eta\,T}
\;+\;
\frac{2\eta L_{\max}}{B}\,
\frac{1}{T}\sum_{t=0}^{T-1}\frac{1}{\sqrt{t+1}} .
\]
In particular, choosing a constant $\eta\le \frac{1}{L_{\max}}$ and using $\sum_{t=0}^{T-1}\frac{1}{\sqrt{t+1}}\le 2\sqrt{T}$ yields
\[
\boxed{\;
\frac{1}{T}\sum_{t=0}^{T-1}\E\!\big[\norm{\nabla f(w_t)}^{2}\big]
\;\le\;
\frac{2B\big(f(w_0)-f^{\star}\big)}{\eta\,T}
\;+\;
\frac{4\eta L_{\max}}{B\sqrt{T}}
\;}
\]
Thus the expected squared norm of the gradient decays as $O\!\left(\frac{1}{\sqrt{T}}\right)$.
\end{theorem}

\section*{Theoretical Properties}
\begin{itemize}
\item \textbf{Stability.} The bound requires only $L_{\max}$-type smoothness; no convexity is needed. Hence the algorithm is stable for a broad class of non‑convex problems.
\item \textbf{Hyperparameter Sensitivity.} The rate contains the product $\eta L_{\max}$. Choosing $\eta$ too large violates the smoothness‑step condition and may cause divergence; a safe choice is $\eta\le 1/L_{\max}$.
\item \textbf{Comparison to Full‑gradient SGD.} When $B=d$ (each coordinate is its own block) the bound recovers the classic $O(1/\sqrt{T})$ rate for stochastic gradient descent with a single sampled component.
\item \textbf{Special Cases.}
    \begin{enumerate}
        \item \emph{Uniform Lipschitz constants} $L_b\equiv L$: $L_{\max}=L$.
        \item \emph{Deterministic cyclic selection} (instead of random) yields the same bound up to a factor $B$ in the numerator.
    \end{enumerate}
\end{itemize}

\section*{Discussion}
The theorem shows that even though only a random block is updated at each iteration, the algorithm makes progress measured by the full‑gradient norm. The $1/\sqrt{T}$ decay matches the optimal rate for first‑order methods on smooth non‑convex objectives, confirming that random block selection does not deteriorate asymptotic performance while reducing per‑iteration computational cost.

\newpage
\section*{Preliminaries (Definitions and Lemmas)}
\begin{lemma}[Block Smoothness Inequality]\label{lem:smooth}
If $f$ satisfies Assumption \ref{ass:smooth}, then for any $w\in\R^{d}$, any block $b$, and any vector $v\in\R^{|\block{b}|}$,
\[
f\big(w-U_{b}v\big)\le f(w)-\langle \nabla_{\block{b}} f(w),v\rangle+\frac{L_b}{2}\|v\|^{2}.
\]
\end{lemma}
\begin{proof}
Apply the definition of $L_b$‑smoothness with $v\gets -v$ and rearrange terms. ∎
\end{proof}

\section*{Main Proof (Step‑by‑Step Derivation)}
We prove Theorem \ref{thm:main}. Throughout the proof, expectations are taken with respect to the random block $i_t$ conditioned on the history $\mathcal{F}_t:=\sigma(w_0,\dots,w_t)$.

\bigskip
\noindent\textbf{[Step 1]} (One‑step descent using Lemma \ref{lem:smooth}).  
For the selected block $i_t$ and the update $w_{t+1}=w_t-\eta_tU_{i_t}\nabla_{\block{i_t}}f(w_t)$, Lemma \ref{lem:smooth} with $v=\eta_t\nabla_{\block{i_t}}f(w_t)$ gives
\begin{equation}\label{eq:descent}
f(w_{t+1})\le f(w_t)-\eta_t\|\nabla_{\block{i_t}}f(w_t)\|^{2}
+\frac{L_{i_t}\eta_t^{2}}{2}\|\nabla_{\block{i_t}}f(w_t)\|^{2}.
\end{equation}
\noindent\textbf{[Step 2]} (Take conditional expectation).  
Conditioning on $\mathcal{F}_t$ and using $\Pr(i_t=b)=1/B$,
\begin{align}
\E\!\big[f(w_{t+1})\mid\mathcal{F}_t\big]
&\le f(w_t)-\eta_t\frac{1}{B}\sum_{b=1}^{B}\|\nabla_{\block{b}}f(w_t)\|^{2}
+\frac{\eta_t^{2}}{2B}\sum_{b=1}^{B}L_{b}\|\nabla_{\block{b}}f(w_t)\|^{2}\nonumber\\
&= f(w_t)-\frac{\eta_t}{B}\sum_{b=1}^{B}\Bigl(1-\frac{L_{b}\eta_t}{2}\Bigr)\|\nabla_{\block{b}}f(w_t)\|^{2}.
\label{eq:exp-descent}
\end{align}
\noindent\textbf{[Step 3]} (Use stepsize bound).  
Assumption \ref{ass:stepsize} guarantees $\eta_t\le 1/L_{\max}\le 2/L_{b}$ for every $b$, hence $1-\frac{L_{b}\eta_t}{2}\ge \frac{1}{2}$. Plugging this into \eqref{eq:exp-descent} yields
\begin{equation}\label{eq:exp-descent2}
\E\!\big[f(w_{t+1})\mid\mathcal{F}_t\big]\le
f(w_t)-\frac{\eta_t}{2B}\sum_{b=1}^{B}\|\nabla_{\block{b}}f(w_t)\|^{2}.
\end{equation}

\noindent\textbf{[Step 4]} (Relate block gradients to full gradient).  
Since the blocks form a partition,
\[
\sum_{b=1}^{B}\|\nabla_{\block{b}}f(w_t)\|^{2}= \|\nabla f(w_t)\|^{2}.
\tag{4.1}
\]
Insert (4.1) into \eqref{eq:exp-descent2}:
\begin{equation}\label{eq:exp-descent3}
\E\!\big[f(w_{t+1})\mid\mathcal{F}_t\big]\le
f(w_t)-\frac{\eta_t}{2B}\,\|\nabla f(w_t)\|^{2}.
\end{equation}

\noindent\textbf{[Step 5]} (Take total expectation).  
Denote $\Delta_t:=\E[f(w_t)]-f^{\star}\ge0$. Taking expectations on both sides of \eqref{eq:exp-descent3} gives
\[
\Delta_{t+1}\le \Delta_t-\frac{\eta_t}{2B}\,\E\!\big[\|\nabla f(w_t)\|^{2}\big].
\tag{5.1}
\]

\noindent\textbf{[Step 6]} (Summation and telescoping).  
Rearrange (5.1) and sum for $t=0,\dots,T-1$:
\[
\sum_{t=0}^{T-1}\frac{\eta_t}{2B}\,\E\!\big[\|\nabla f(w_t)\|^{2}\big]
\le \Delta_0-\Delta_T\le f(w_0)-f^{\star}.
\tag{6.1}
\]

\noindent\textbf{[Step 7]} (Insert the explicit stepsize $\eta_t=\eta/\sqrt{t+1}$).  
Equation (6.1) becomes
\[
\frac{\eta}{2B}\sum_{t=0}^{T-1}\frac{1}{\sqrt{t+1}}\,
\E\!\big[\|\nabla f(w_t)\|^{2}\big]
\le f(w_0)-f^{\star}.
\tag{7.1}
\]

\noindent\textbf{[Step 8]} (Upper bound the harmonic‑type sum).  
Since $1/\sqrt{t+1}\ge 1/\sqrt{T}$ for all $t\le T-1$, we have
\[
\sum_{t=0}^{T-1}\frac{1}{\sqrt{t+1}}\ge \frac{T}{\sqrt{T}}= \sqrt{T}.
\tag{8.1}
\]
Using (8.1) in (7.1) gives a crude bound
\[
\frac{\eta\sqrt{T}}{2B}\,\frac{1}{T}\sum_{t=0}^{T-1}\E\!\big[\|\nabla f(w_t)\|^{2}\big]
\le f(w_0)-f^{\star},
\]
which after rearranging yields the first term of the theorem:
\[
\frac{1}{T}\sum_{t=0}^{T-1}\E\!\big[\|\nabla f(w_t)\|^{2}\big]
\le \frac{2B\big(f(w_0)-f^{\star}\big)}{\eta\,\sqrt{T}}.
\tag{8.2}
\]

\noindent\textbf{[Step 9]} (Refine the bound keeping the exact sum).  
From (6.1) with the explicit $\eta_t$ we obtain
\[
\frac{1}{T}\sum_{t=0}^{T-1}\E\!\big[\|\nabla f(w_t)\|^{2}\big]
\le
\frac{2B\big(f(w_0)-f^{\star}\big)}{\eta\,T}
\;+\;
\frac{2\eta L_{\max}}{B}\,
\frac{1}{T}\sum_{t=0}^{T-1}\frac{1}{\sqrt{t+1}} .
\tag{9.1}
\]
The second term arises from the slack we introduced in Step 3 when we replaced $1-\frac{L_b\eta_t}{2}$ by $\frac12$. Keeping the exact factor yields precisely the additive term displayed in (9.1). This is the statement of Theorem \ref{thm:main}.

\noindent\textbf{[Step 10]} (Simplify using $\sum_{t=0}^{T-1}1/\sqrt{t+1}\le 2\sqrt{T}$).  
Applying the elementary bound $\sum_{k=1}^{T}k^{-1/2}\le 2\sqrt{T}$ to (9.1) gives the compact rate
\[
\frac{1}{T}\sum_{t=0}^{T-1}\E\!\big[\|\nabla f(w_t)\|^{2}\big]
\le
\frac{2B\big(f(w_0)-f^{\star}\big)}{\eta\,T}
\;+\;
\frac{4\eta L_{\max}}{B\sqrt{T}} .
\tag{10.1}
\]
This completes the proof. ∎

\section*{Rate Derivation (With Constants)}
From (10.1) we identify two competing terms:
\[
\underbrace{\frac{2B\big(f(w_0)-f^{\star}\big)}_{\displaystyle C_1}
\frac{1}{\eta T}
\quad\text{and}\quad
\underbrace{\frac{4L_{\max}}{B}}_{\displaystyle C_2}
\frac{\eta}{\sqrt{T}} .
\]
Choosing $\eta = \sqrt{\frac{C_1}{C_2}}\;T^{-1/4}$ balances the two contributions and yields the optimal rate $O(T^{-1/2})$. For a fixed stepsize $\eta\le 1/L_{\max}$ the bound simplifies to $O(1/\sqrt{T})$ as stated.

\section*{Special Cases}
\begin{itemize}
\item \textbf{Uniform Lipschitz constants.} If $L_b\equiv L$, then $L_{\max}=L$ and the constants $C_1,C_2$ become $2B\Delta_0$ and $4L/B$, respectively.
\item \textbf{Full‑gradient case ($B=1$).} The algorithm reduces to ordinary gradient descent and the bound recovers the classic $O(1/\sqrt{T})$ rate for non‑convex smooth optimization.
\item \textbf{Large number of blocks.} As $B$ grows, the first term scales linearly with $B$ while the second term scales inversely with $B$, reflecting the trade‑off between per‑iteration progress and variance reduction.
\end{itemize}

\end{document}